\chapter{记号表}

本附录按主题汇总本书中反复出现的记号。
“首次出现”一栏给出主要章节或附录，
便于读者在查阅记号时回到原始语境。
同一记号在不同章节偶有轻微语义变化时，
应以所在章节上下文为准。

\section{集合论记号}

\begin{longtable}{p{0.22\textwidth}p{0.46\textwidth}p{0.22\textwidth}}
\toprule
记号 & 含义 & 首次出现 \\
\midrule
\endfirsthead
\toprule
记号 & 含义 & 首次出现 \\
\midrule
\endhead
$\varnothing$ & 空集 & 第 1 章 \\
$x\in A$ & 元素 $x$ 属于集合 $A$ & 第 1 章 \\
$A\subseteq B$ & $A$ 是 $B$ 的子集 & 第 1 章 \\
$A\subsetneq B$ & 真子集 & 第 1 章 \\
$A\cup B$ & 并集 & 第 1 章 \\
$A\cap B$ & 交集 & 第 1 章 \\
$A\setminus B$ & 差集 & 第 1 章 \\
$A\times B$ & 笛卡儿积 & 第 1 章，附录 A \\
$\#A$ & 有限集合 $A$ 的基数 & 第 1 章 \\
$f:A\to B$ & 从 $A$ 到 $B$ 的映射 & 第 1 章 \\
$f^{-1}(S)$ & 子集 $S\subseteq B$ 的逆像 & 第 1 章，附录 A \\
$\mathcal{P}(A)$ & $A$ 的幂集 & 第 1 章 \\
\bottomrule
\end{longtable}

\section{群论记号}

\begin{longtable}{p{0.22\textwidth}p{0.46\textwidth}p{0.22\textwidth}}
\toprule
记号 & 含义 & 首次出现 \\
\midrule
\endfirsthead
\toprule
记号 & 含义 & 首次出现 \\
\midrule
\endhead
$G,H,K$ & 群 & 第 1 章 \\
$e$ 或 $1$ & 群单位元 & 第 1 章 \\
$g^{-1}$ & 元素 $g$ 的逆元 & 第 1 章 \\
$\ord(g)$ & 元素 $g$ 的阶 & 第 1 章 \\
$\langle S\rangle$ & 由子集 $S$ 生成的子群 & 第 1 章 \\
$H\le G$ & $H$ 是 $G$ 的子群 & 第 1 章 \\
$H\triangleleft G$ & $H$ 是 $G$ 的正规子群 & 第 1 章 \\
$G/H$ & 商群 & 第 1 章 \\
$\Aut(G)$ & 群 $G$ 的自同构群 & 第 2 章 \\
$\Inn(G)$ & 群 $G$ 的内自同构群 & 第 2 章 \\
$\Syl_p(G)$ & $G$ 的 Sylow $p$-子群集合 & 第 2 章 \\
$\Stab_G(x)$ & 群作用下点 $x$ 的稳定子 & 第 3 章 \\
$\Orb_G(x)$ 或 $\Orb(x)$ & 点 $x$ 的轨道 & 第 3 章 \\
$\Fix_G(X)$ 或 $\Fix(X)$ & 被整个群固定的点集 & 第 3 章 \\
$\GL_n(F)$ & 一般线性群 & 第 3 章 \\
$\SL_n(F)$ & 特殊线性群 & 第 3 章 \\
$\PSL_2(\Z)$ & $\SL_2(\Z)$ 模去中心后的射影群 & 第 3 章 \\
\bottomrule
\end{longtable}

\section{环论记号}

\begin{longtable}{p{0.22\textwidth}p{0.46\textwidth}p{0.22\textwidth}}
\toprule
记号 & 含义 & 首次出现 \\
\midrule
\endfirsthead
\toprule
记号 & 含义 & 首次出现 \\
\midrule
\endhead
$R,S,T$ & 环 & 第 5 章 \\
$1_R$ & 环 $R$ 的乘法单位元 & 第 5 章 \\
$R^{\times}$ & 环 $R$ 的单位群 & 第 5 章 \\
$R[x]$ & 以 $R$ 为系数的多项式环 & 第 5 章 \\
$I\lhd R$ & 双边理想 & 第 5 章 \\
$R/I$ & 商环 & 第 5 章 \\
$\Ann(M)$ & 模 $M$ 的消去子 & 第 7 章 \\
$\Spec(R)$ & 素理想谱 & 第 6 章 \\
$\Rad(I)$ & 理想 $I$ 的根基 & 第 6 章 \\
$\Char(R)$ & 环 $R$ 的特征 & 第 5 章 \\
$\mathrm{Frac}(R)$ & 整环 $R$ 的分式域 & 第 6 章 \\
$\mathcal{H}_k$ & 作用在权 $k$ 模形式空间上的 Hecke 代数 & 第 10 章 \\
$T_n$ & 第 $n$ 个 Hecke 算子 & 第 10 章 \\
\bottomrule
\end{longtable}

\section{模论记号}

\begin{longtable}{p{0.22\textwidth}p{0.46\textwidth}p{0.22\textwidth}}
\toprule
记号 & 含义 & 首次出现 \\
\midrule
\endfirsthead
\toprule
记号 & 含义 & 首次出现 \\
\midrule
\endhead
$M,N,P$ & 左 $R$-模 & 第 7 章 \\
$\Hom_R(M,N)$ & $R$-模同态集合 & 第 7 章 \\
$\End_R(M)$ & 模 $M$ 的自同态环 & 第 7 章 \\
$\Ker(f)$ & 同态 $f$ 的核 & 第 7 章 \\
$\Ima(f)$ 或 $\im(f)$ & 同态 $f$ 的像 & 第 7 章 \\
$\coker(f)$ & 同态 $f$ 的余核 & 附录 A \\
$\rank(M)$ & 模或线性映射的秩 & 第 7 章 \\
$\ses{A}{B}{C}$ & 短正合列 $0\to A\to B\to C\to 0$ & 第 7 章 \\
$M\otimes_R N$ & 张量积 & 第 7 章，附录 A \\
$M^*=\Hom_R(M,R)$ & 对偶模 & 第 7 章 \\
$P$ 是投射的 & 提升性质成立，等价于 $\Hom_R(P,-)$ 正合 & 第 7 章，附录 A \\
$M$ 是平坦的 & $-\otimes_R M$ 正合 & 第 7 章，附录 A \\
\bottomrule
\end{longtable}

\section{域论记号}

\begin{longtable}{p{0.22\textwidth}p{0.46\textwidth}p{0.22\textwidth}}
\toprule
记号 & 含义 & 首次出现 \\
\midrule
\endfirsthead
\toprule
记号 & 含义 & 首次出现 \\
\midrule
\endhead
$F\subseteq K\subseteq L$ & 域扩张塔 & 第 8 章 \\
$[K:F]$ & 域扩张次数 & 第 8 章 \\
$\Aut_F(K)$ & 固定 $F$ 的自同构群 & 第 8 章 \\
$\Gal(L/F)$ & Galois 群 & 第 9 章 \\
$\mathrm{Fix}(H)$ & 子群 $H$ 的不动域 & 第 9 章 \\
$\mathrm{Gal}(L/E)$ & 中间域 $E$ 对应的子群 & 第 9 章 \\
$\Tr_{L/F}$ & 域迹 & 第 8 章 \\
$\Nm_{L/F}$ & 域范数 & 第 8 章 \\
$\Frob_p$ & Frobenius 元或 Frobenius 共轭类 & 第 9 章，第 10 章 \\
$\overline{\Q}$ & $\Q$ 的代数闭包 & 第 9 章 \\
$\overline{\Q}_\ell$ & $\ell$-进代数闭包 & 第 10 章 \\
\bottomrule
\end{longtable}

\section{线性代数记号}

\begin{longtable}{p{0.22\textwidth}p{0.46\textwidth}p{0.22\textwidth}}
\toprule
记号 & 含义 & 首次出现 \\
\midrule
\endfirsthead
\toprule
记号 & 含义 & 首次出现 \\
\midrule
\endhead
$\Mat_{m\times n}(F)$ & $m\times n$ 矩阵全体 & 第 3 章 \\
$\mat{a&b\\c&d}$ & 用于排版矩阵的简写命令 & 导言区 \\
$\dmat{a&b\\c&d}$ & 用于排版行列式的简写命令 & 导言区 \\
$\diag(a_1,\dots,a_n)$ & 对角矩阵 & 第 3 章 \\
$\Tr(A)$ & 矩阵或线性变换的迹 & 第 4 章 \\
$\det(A)$ & 行列式 & 第 3 章 \\
$\inner{v}{w}$ & 内积或配对 & 第 4 章 \\
$V^*$ & 向量空间或模的对偶空间 & 第 4 章，附录 A \\
$V^{**}$ & 双对偶空间 & 附录 A \\
$\lambda$ & 特征值或标量参数 & 第 3 章起 \\
$\chi_V$ & 表示或线性变换的特征标 & 第 4 章 \\
\bottomrule
\end{longtable}

\section{模形式记号}

\begin{longtable}{p{0.22\textwidth}p{0.46\textwidth}p{0.22\textwidth}}
\toprule
记号 & 含义 & 首次出现 \\
\midrule
\endfirsthead
\toprule
记号 & 含义 & 首次出现 \\
\midrule
\endhead
$\HH$ & 上半平面 $\{z\in\C:\Im z>0\}$ & 第 10 章 \\
$q=e^{2\pi i z}$ & Fourier 展开的局部参数 & 第 10 章 \\
$\Mk_k(\Gamma)$ & 权 $k$、群 $\Gamma$ 上的模形式空间 & 第 10 章 \\
$\Sk_k(\Gamma)$ & 权 $k$、群 $\Gamma$ 上的尖点形式空间 & 第 10 章 \\
$f\mid_k \gamma$ & 权 $k$ 的 slash 作用 & 第 10 章 \\
$E_4,E_6$ & 经典 Eisenstein 级数 & 第 10 章 \\
$\Delta$ & 判别尖点形式 & 第 10 章 \\
$\tau(n)$ & Ramanujan $\tau$ 函数 & 第 10 章 \\
$L(f,s)$ & 模形式 $f$ 的 $L$-函数 & 第 10 章 \\
$\GammaO(N)$ & 零类同余子群 $\Gamma_0(N)$ & 第 10 章 \\
$\GammaI(N)$ & 第一类同余子群 $\Gamma_1(N)$ & 第 10 章 \\
$\Gamma(N)$ & 主同余子群 & 第 10 章 \\
$a_n(f)$ 或 $a_n$ & 模形式 Fourier 系数，也即特征形式时的 Hecke 特征值 & 第 10 章 \\
$newform$ & 经过 Atkin--Lehner 理论筛选的归一化新形式 & 第 10 章 \\
\bottomrule
\end{longtable}

\section{范畴论记号}

\begin{longtable}{p{0.22\textwidth}p{0.46\textwidth}p{0.22\textwidth}}
\toprule
记号 & 含义 & 首次出现 \\
\midrule
\endfirsthead
\toprule
记号 & 含义 & 首次出现 \\
\midrule
\endhead
$\mathcal{C},\mathcal{D}$ & 范畴 & 附录 A \\
$\Hom_{\mathcal{C}}(A,B)$ & 范畴 $\mathcal{C}$ 中从 $A$ 到 $B$ 的态射集 & 附录 A \\
$\Id_A$ & 对象 $A$ 的恒等态射 & 附录 A \\
$\mathcal{C}^{\mathrm{op}}$ & 反范畴 & 附录 A \\
$F:\mathcal{C}\to\mathcal{D}$ & 协变函子 & 附录 A \\
$F:\mathcal{C}^{\mathrm{op}}\to\mathcal{D}$ & 逆变函子 & 附录 A \\
$\eta:F\Rightarrow G$ & 从函子 $F$ 到 $G$ 的自然变换 & 附录 A \\
$\mathrm{Nat}(F,G)$ & 自然变换集合 & 附录 A \\
$F\dashv U$ & $F$ 左伴随于 $U$ & 附录 A \\
$\prod_i A_i$ & 积 & 附录 A \\
$\coprod_i A_i$ & 余积 & 附录 A \\
$\Hom(A,-)$ & 可表函子的典型形式 & 附录 A \\
$h^A$ & 常用记号，表示 $\Hom(A,-)$ & 附录 A \\
\bottomrule
\end{longtable}
